\documentclass[11pt,a4paper]{article}
\usepackage[T1]{fontenc}
\usepackage{lmodern}
\usepackage[margin=27mm]{geometry}
\usepackage{amsmath,amssymb,amsthm,mathtools}
\usepackage{booktabs,microtype}
\usepackage[hidelinks]{hyperref}
\usepackage{url}
\hypersetup{
 pdftitle={Singular energy minima and fixed-product structure for all-odd complex unit gain dumbbells},
 pdfauthor={Xiaojian Zeng}
}
\newtheorem{theorem}{Theorem}
\newtheorem{proposition}[theorem]{Proposition}
\newtheorem{lemma}[theorem]{Lemma}
\theoremstyle{remark}
\newtheorem{remark}[theorem]{Remark}
\DeclareMathOperator{\Rea}{Re}
\DeclareMathOperator{\sgn}{sgn}
\DeclareMathOperator{\tr}{tr}
\DeclareMathOperator{\arsinh}{arsinh}
\newcommand{\dd}{\,dt}
\newcommand{\E}{\mathcal E}
\newcommand{\R}{\mathbb R}
\newcommand{\T}{\mathbb T}
\newcommand{\eps}{\varepsilon}
\title{Singular energy minima and fixed-product structure\\
for all-odd complex unit gain dumbbells}
\author{Xiaojian Zeng\\
\small PBC School of Finance, Tsinghua University, 100190\\
\small \href{mailto:xiaojianzeng@foxmail.com}{xiaojianzeng@foxmail.com}}
\date{}
\begin{document}
\maketitle
\begin{abstract}
We study energy minimization on a fixed dumbbell graph whose two cycle
lengths and connecting-path length are odd. The full minimum-energy and
switching-equality classification remains unresolved. We give an explicit
expansion at every interior parameter pair with a zero eigenvalue, replacing
the unqualified smooth stationarity equations of Huang and Pereyra by a
tangential equation and a normal inequality. A strict normal inequality
also gives a strict local minimum. On every positive fixed-product curve,
the integrated energy is strictly convex in the logarithmic gain ratio.
A length-dependent comparison forces its minimizer to the boundary in an
explicit range. Exact examples distinguish the possibilities: the two
half-diagonal points are the only global minima for two triangles joined
by one edge, whereas the analogous points for a five-edge connector admit
a certified direction of descent.
\end{abstract}

\section{The problem and the contribution}

A complex unit gain graph assigns $\varphi_{uv}\in\T=\{z\in\mathbb C:|z|=1\}$
to each oriented edge, with $\varphi_{vu}=\overline{\varphi_{uv}}$.
Its adjacency matrix $A_\varphi$ is Hermitian, and its energy is
\[
 \E(\varphi)=\sum_{j=1}^{N}|\lambda_j(A_\varphi)|
            =\tr|A_\varphi|.
\]
We use the gain-graph and switching conventions of Reff~\cite{Reff2012}.
Thus switching is diagonal unitary conjugation. Energy is the Schatten
$1$-norm of the adjacency matrix. Matrix absolute values also enter
Schatten-norm comparison problems. Zeng, Liu and Ratnavelu
\cite{ZengLiuRatnavelu2026} certify a Schatten-norm counterexample at exponent
$3/2$ by exact rational inequalities. Here the objective
is the trace norm on a prescribed gain-graph family, and the small
examples below likewise use exact sign certificates.

Let $D=D_{r,s,\ell}$ consist of disjoint cycles $C_r,C_s$ whose designated
vertices are joined by an $\ell$-edge path with new internal vertices.
Throughout,
\[
 r,s\geq3,\quad \ell\geq1\quad\hbox{are odd},\qquad
 N=r+s+\ell-1,\quad L=\frac{\ell-1}{2}.
\]
Fix orientations on the two cycles, with gain products $\gamma_r,\gamma_s$.
Deleting one edge from each cycle gives a spanning tree. Its gains can
be switched to $1$, and the two cycle products then determine the
switching class. Write
\begin{equation}\label{eq:normalization}
 a=\Rea\gamma_r,\quad b=\Rea\gamma_s,\qquad
 \alpha=(-1)^{(r+1)/2}a,\quad
 \beta=(-1)^{(s+1)/2}b .
\end{equation}
Every $(\alpha,\beta)\in[-1,1]^2$ is attainable. These two real parameters
determine the spectrum but do not distinguish all switching classes:
each pair $(a,b)$ lifts to
$\gamma_r=a\pm i\sqrt{1-a^2}$ and
$\gamma_s=b\pm i\sqrt{1-b^2}$, counting coincident endpoint choices once.
No quotient by graph automorphisms is taken.

Samanta and Kannan~\cite[Theorem 4.3 and Section 5]{SamantaKannan2025}
give a matching-polynomial expansion and extremal energy comparisons
for unicyclic gain graphs. Huang and Pereyra~\cite{HuangPereyra2026}
study the dumbbell family and leave the all-odd minimum problem
undetermined. Their Theorem 4.6 minimizes the Coulson kernel at each
fixed integration variable. Their Theorem 4.7 then states smooth
stationarity equations for every interior energy minimizer. However,
the energy may have a simple zero-eigenvalue cusp. Their Proposition 4.9
already gives the half-diagonal energy for $D_{3,3,1}$, and their
discussion of Figure 2 already recognizes this cusp.

Our contribution is an explicit correction at singular parameters,
including a strict local sufficiency test, together with uniqueness for
the \emph{integrated} fixed-product problem and a length-based boundary
criterion. We prove global minimality in the $D_{3,3,1}$ test case and
give an exact contrasting derivative for $D_{3,3,5}$. The known
half-diagonal energy and the general nonsmooth spectral mechanism are
not contributions of this note. Differentiation and generalized
gradients of spectral functions are established theory
\cite{Lewis1996}; the argument below isolates the zero root explicitly.
The comparison concerns version 2 of~\cite{HuangPereyra2026}.
No equivalent graph-specific correction was found in the primary
sources inspected through 14 September 2026; this is a bounded
literature assessment.

\section{Kernel and statements}

For a graph $G$, let $M_k(G)$ count its $k$-edge matchings and put
\[
 m_G(x)=\sum_{k\geq0}(-1)^k M_k(G)x^{|V(G)|-2k}.
\]
The path $P_j$ has $j$ vertices; $P_0$ is empty, with $m_{P_0}=1$.
The characteristic polynomial $P(x;\alpha,\beta)=\det(xI-A_\varphi)$,
with $a,b$ recovered from~\eqref{eq:normalization}, is
\begin{equation}\label{eq:matching}
 P(x;\alpha,\beta)
 =m_D(x)-2a\,m_{D-C_r}(x)-2b\,m_{D-C_s}(x)
       +4ab\,m_{P_{\ell-1}}(x).
\end{equation}
This is~\cite[Theorem 2.10]{HuangPereyra2026}, also a specialization
of~\cite[Theorem 4.3]{SamantaKannan2025}. For completeness, a determinant
permutation contributes edge transpositions and oriented simple cycles.
Each edge transposition has weight $-1$, and the two orientations of
a cycle together have weight $-2$ times its real cycle gain. The only
nonempty unions of disjoint cycles in $D$ are $C_r,C_s,C_r\cup C_s$;
summing the remaining matchings gives~\eqref{eq:matching}.

We specify the positive polynomials in the kernel so that the formulas
are directly evaluable. With an indeterminate $v$, define
\begin{align*}
 f_{-1}(v)&=0,\quad f_0(v)=1,\quad
 f_j(v)=f_{j-1}(v)+vf_{j-2}(v)\quad(j\geq1),\\
 g_j(v)&=f_j(v)+vf_{j-2}(v)\quad(j\geq3).
\end{align*}
These are the matching generating polynomials for paths and cycles.
Set
\begin{align*}
 M(v)&=\sum_{k\geq0}M_k(D)v^k,& T(v)&=f_{\ell-1}(v),\\
 R(v)&=f_{\ell-1}(v)g_s(v)+vf_{\ell-2}(v)f_{s-1}(v),\\
 S(v)&=f_{\ell-1}(v)g_r(v)+vf_{\ell-2}(v)f_{r-1}(v).
\end{align*}
The last two identities follow by separating matchings according to
whether the tail uses its edge into the remaining cycle; they hold
also for $\ell=1$. For $t>0$, put
\begin{equation}\label{eq:polynomials}
 Z=M(t^2),\quad W=4t^{r+s}T(t^2),\quad
 X=2t^rR(t^2),\quad Y=2t^sS(t^2).
\end{equation}
All four quantities are strictly positive. In the notation
\begin{equation}\label{eq:kernel}
 A=Z-\alpha\beta W,\qquad B=\alpha X+\beta Y,\qquad C=A^2+B^2,
\end{equation}
Lemma 4.2 of~\cite{HuangPereyra2026} gives
\begin{equation}\label{eq:energy-integral}
 C=|t^NP(i/t;\alpha,\beta)|^2,\qquad
 \E(\alpha,\beta)=\frac1\pi\int_0^\infty\frac{\log C}{t^2}\dd .
\end{equation}
Indeed, multiplying~\eqref{eq:matching} by $t^N$ and using
$t^j m_G(i/t)=i^j\sum_k M_k(G)t^{2k}$ for a $j$-vertex graph gives
$t^NP(i/t)=i^N(A-iB)$. The integral is the real-rooted case of the
polynomial Coulson formula~\cite[Theorem 1]{MateljevicBozinGutman2010}.
It follows directly here by summing the elementary identity
\begin{equation}\label{eq:scalar-integral}
 \frac1\pi\int_0^\infty
       \frac{\log(1+t^2z^2)}{t^2}\dd=|z|,\qquad z\in\R,
\end{equation}
which follows by integration by parts, including $z=0$.

Define, at the parameter pair in question,
\begin{align}
 I_\alpha&=\int_0^\infty\frac{-\beta WA+XB}{t^2C}\dd,&
 I_\beta&=\int_0^\infty\frac{-\alpha WA+YB}{t^2C}\dd,\label{eq:integrals}\\
 H&=(s+L)\alpha+(r+L)\beta,&
 J&=\alpha I_\alpha+\beta I_\beta.\label{eq:HJ}
\end{align}
Convergence at the singular points is justified in Section~\ref{sec:singular}.

\begin{theorem}[Singular expansion and optimality]\label{thm:singular}
Write $\eps=(-1)^{N/2}$. On the closed parameter square,
\begin{equation}\label{eq:coefficients}
 P(0;\alpha,\beta)=\eps(1-4\alpha\beta),\qquad
 P_x(0;\alpha,\beta)=-2\eps H.
\end{equation}
The zero-eigenvalue locus is $\alpha\beta=1/4$, and its zero eigenvalue
is simple everywhere, including its boundary points.
For each interior point $q=(\alpha,\beta)$ on this locus,
\begin{equation}\label{eq:expansion}
 \begin{split}
 \E(q+h)=\E(q)&+\frac2\pi(I_\alpha h_\alpha+I_\beta h_\beta)\\
 &+\frac{2|\beta h_\alpha+\alpha h_\beta|}{|H|}
       +O(\|h\|^2).
 \end{split}
\end{equation}
Thus necessary conditions for an interior singular local minimum are
\begin{equation}\label{eq:conditions}
 \alpha I_\alpha=\beta I_\beta,\qquad
 |J|\leq\frac{\pi}{2|H|}.
\end{equation}
They are equivalent to nonnegativity of all first directional
derivatives. If the equality and the \emph{strict} inequality
in~\eqref{eq:conditions} hold, $q$ is a strict local minimum.
The case of equality in the normal bound is not classified.
At an interior point off $\alpha\beta=1/4$, the usual necessary
conditions $I_\alpha=I_\beta=0$ are valid.
\end{theorem}

\begin{theorem}[Fixed-product uniqueness and boundary criterion]\label{thm:product}
For each $0<p<1$, the energy on
$\{(\alpha,\beta)\in(0,1]^2:\alpha\beta=p\}$ is strictly convex in
\[
 u=\frac12\log(\alpha/\beta),\qquad
 \frac12\log p\leq u\leq-\frac12\log p .
\]
It has exactly one constrained minimizer. At $p=1$ the curve is a
singleton. If $r=s$, the minimizer is $(\sqrt p,\sqrt p)$.
If $r<s$ and $p<1$, its coordinates satisfy $\alpha<\beta$, and
\begin{equation}\label{eq:ratio-bound}
 \frac{X(t)}{Y(t)}>\frac{s+L}{r+L}\qquad(t>0).
\end{equation}
In particular, if $p\geq(r+L)/(s+L)$, the unique constrained
minimizer is $(p,1)$. Hence
\begin{equation}\label{eq:length-bound}
 s+L\geq4(r+L)
\end{equation}
forces the positive singular-curve minimum to $(1/4,1)$ and excludes
every interior singular local minimum of the full-square energy.
Cycle exchange gives the corresponding statements for $s<r$,
and simultaneous negation gives the negative-branch statements.
No full-square local or global minimality of these boundary
endpoints is asserted.
\end{theorem}

\section{Proof of the fixed-product structure}

Write $\alpha=\sqrt p\,e^u$, $\beta=\sqrt p\,e^{-u}$. At fixed
$p>0$ and $t>0$, the kernel takes the form
\[
 C=d+z+w,\qquad
 d=(Z-pW)^2+2pXY>0,\quad
 z=pX^2e^{2u}>0,\quad w=pY^2e^{-2u}>0.
\]
A direct calculation gives
\begin{equation}\label{eq:strict-convex}
 (\log C)_{uu}=\frac{4\{d(z+w)+4zw\}}{(d+z+w)^2}>0.
\end{equation}
Integrating the strict convexity inequality in~\eqref{eq:energy-integral}
proves strict convexity of $\E$ without exchanging differentiation and
integration. Continuity and compactness give existence, and strict
convexity gives uniqueness. The first two $u$ derivatives can also be
integrated: before division by $t^2$ they are
$O(t^{2\min(r,s)})$ as $t\downarrow0$, and they are uniformly bounded
as $t\to\infty$ on the fixed compact $u$ interval. Symmetry proves the
diagonal assertion when $r=s$.

To compare unequal cycle lengths, set $\eta=\arsinh(1/(2t))>0$.
The characteristic roots of the path recurrence are $te^\eta$
and $-te^{-\eta}$. For odd $k$ they give
\[
 g_k(t^2)=2t^k\sinh(k\eta),\qquad
 f_{k-1}(t^2)=t^{k-1}\frac{\cosh(k\eta)}{\cosh\eta}.
\]
Also
$f_{\ell-2}(t^2)=t^{\ell-2}\sinh((\ell-1)\eta)/\cosh\eta$,
which agrees with $f_{-1}=0$ when $\ell=1$.
Substituting in~\eqref{eq:polynomials} yields
\begin{equation}\label{eq:hyperbolic}
 \frac XY=
 \frac{\sinh(s\eta)+c_\eta\cosh(s\eta)}
      {\sinh(r\eta)+c_\eta\cosh(r\eta)},\qquad
 c_\eta=\frac{\sinh((\ell-1)\eta)}
              {2\cosh(\ell\eta)\cosh\eta}.
\end{equation}
Since $\cosh(\ell\eta)\cosh\eta\geq\cosh((\ell-1)\eta)$ and
$\tanh y\leq y$ for $y\geq0$, we have $0\leq c_\eta\leq L\eta$.
For fixed $c\geq0$, the function
\[
 z\longmapsto\frac{\sinh z+c\cosh z}{z+c}\qquad(z>0)
\]
is strictly increasing: its derivative numerator is
$z\cosh z-\sinh z+cz\sinh z+c^2\sinh z>0$.
For $r<s$, it follows that
\[
 \frac XY>\frac{s+c_\eta/\eta}{r+c_\eta/\eta}
            \geq\frac{s+L}{r+L},
\]
proving~\eqref{eq:ratio-bound}. Moreover,
\begin{equation}\label{eq:u-derivative}
 (\log C)_u=\frac{2(\alpha^2X^2-\beta^2Y^2)}{C}.
\end{equation}
At $u=0$ this is positive, so the constrained minimizer has $u<0$
for $p<1$. If $p\geq(r+L)/(s+L)$, then at the left endpoint
$(\alpha,\beta)=(p,1)$ we have $pX>Y$ for every $t>0$.
The energy derivative there is positive. Strict convexity places
the unique minimum at that endpoint. Apply this with $p=1/4$
to obtain~\eqref{eq:length-bound}, including equality.
Every other point of that singular branch admits constrained descent.
Finally, $\E(-\alpha,-\beta)=\E(\alpha,\beta)$ follows directly
from~\eqref{eq:kernel}, completing the proof of Theorem~\ref{thm:product}.

\section{The simple zero and its directional contribution}\label{sec:singular}

The graph $D$ has exactly one perfect matching: each odd cycle must
use its incident connector edge, and all residual even paths have
forced alternating matchings. The graph $D-C_r$ is $C_s$ with a tail
of $2L$ new vertices. It has exactly $s+L$ matchings leaving one
vertex uncovered. For $L=0$ the count is $s$. Increasing the tail by
two vertices either includes its last edge, recovering the previous
near-perfect count, or leaves the free endpoint uncovered, after
which the remaining graph has one forced perfect matching. Thus the
count increases by one. Similarly $D-C_s$ has $r+L$ such matchings.
Consequently,
\[
 [x]m_{D-C_r}=(-1)^{(N-r-1)/2}(s+L),\qquad
 [x]m_{D-C_s}=(-1)^{(N-s-1)/2}(r+L).
\]
The polynomials $m_D$ and $m_{P_{\ell-1}}$ are even.
Substitution in~\eqref{eq:matching}, with the signs
in~\eqref{eq:normalization}, proves~\eqref{eq:coefficients}.
On $\alpha\beta=1/4$, the two coordinates have the same nonzero
sign, so $H\ne0$. This proves simplicity on the entire locus.

Fix an interior singular point $q$. The implicit function theorem
gives a real analytic root $\lambda$ near $q$, with
\begin{equation}\label{eq:zero-slope}
 \lambda(q)=0,\qquad
 d\lambda(q)[h]=-\frac{2(\beta h_\alpha+\alpha h_\beta)}{H}.
\end{equation}
The other roots stay separated from zero. Their positive and negative
group sums are analytic, even if roots within a group coincide:
each group sum is a contour integral of $xP_x/P$ enclosing that group.
Their absolute sum is therefore an analytic function $F$, and
\[
 \E=F+|\lambda|.
\]

It remains to identify the derivative of $F$ with the displayed
integrals. Factor $P=(x-\lambda)Q$, where $Q$ is monic and real.
Then exactly
\begin{equation}\label{eq:root-removed}
 |t^NP(i/t)|^2=(1+t^2\lambda^2)|t^{N-1}Q(i/t)|^2.
\end{equation}
By~\eqref{eq:scalar-integral}, the first factor contributes $|\lambda|$.
The integral for $Q$ can be differentiated: its differentiated
integrand is uniformly $O(1)$ near $t=0$ and $O(t^{-2})$ at infinity
in a small parameter neighborhood. Monicity and reality of $Q$
give the first estimate, and $Q(0)$ staying away from zero gives
the second. On intermediate compact $t$ intervals, real-rootedness
ensures $Q(i/t)\ne0$.
At $q$ the fixed-$t$ parameter derivative of
$\log(1+t^2\lambda^2)$ is zero. Hence
\begin{equation}\label{eq:smooth-gradient}
 \nabla F(q)=\frac2\pi(I_\alpha,I_\beta).
\end{equation}
This proves convergence in~\eqref{eq:integrals} at $q$.
Taylor expansion of $F$ and $\lambda$, and the Lipschitz property
of the absolute value, prove~\eqref{eq:expansion}. The remainder is
uniform as $h\to0$ for each fixed $q$ and each fixed length triple;
no bound uniform over all graphs is intended. Off the zero locus,
the same domination argument applies directly to $P$.

For the optimality conditions, put $g=\nabla F(q)$ and
$n=(\beta,\alpha)$. All directional derivatives are nonnegative
if and only if
\[
 g=kn\quad\hbox{for some }k\in\R,\qquad |k|\leq2/|H|.
\]
Indeed, the two tangent directions force $g$ to be normal, and
the two normal directions give the bound. Since $\alpha\beta=1/4$,
this is precisely~\eqref{eq:conditions}.

To prove strict sufficiency, use local coordinates $p=\alpha\beta$,
$u=\tfrac12\log(\alpha/\beta)$ on the fixed common-sign branch.
The simple zero vanishes identically on $p=1/4$, so it factors as
\[
 \lambda(p,u)=(p-1/4)a_0(p,u),\qquad a_0(q)=-2/H(q).
\]
Thus $\E(p,u)=F(p,u)+|p-1/4|c(p,u)$, with $c=|a_0|>0$ locally.
Equations~\eqref{eq:conditions} and~\eqref{eq:smooth-gradient} give
\[
 F_u(q)=0,\qquad F_p(q)=4J/\pi,\qquad c(q)=2/|H(q)|.
\]
The strict normal inequality is $|F_p(q)|<c(q)$.
By continuity there is a neighborhood and $\delta>0$ in which
$c(p,u)-|F_p(p',u)|\geq\delta$ for all relevant $p,p',u$.
The mean value theorem yields
\[
 \E(p,u)-\E(1/4,u)\geq\delta|p-1/4|.
\]
By Theorem~\ref{thm:product}, tangential stationarity at $q$
is the unique strict constrained minimum on $p=1/4$.
Together these facts prove strict local minimality and finish
Theorem~\ref{thm:singular}.

\begin{remark}
The tangential equation survives unchanged from
\cite[equation (118)]{HuangPereyra2026}, since
\[
 \alpha I_\alpha-\beta I_\beta
 =\int_0^\infty\frac{\alpha^2X^2-\beta^2Y^2}{t^2C}\dd .
\]
At singular minima the zero identity for $J$ in their equation (119)
is replaced by the interval bound in~\eqref{eq:conditions}.
\end{remark}

\section{Two exact tests}

\begin{proposition}\label{prop:331}
For $D_{3,3,1}$ the full-square global minima occur exactly at
$(\alpha,\beta)=(1/2,1/2)$ and $(-1/2,-1/2)$.
There are eight oriented cycle-product switching classes attaining
the common minimum $2+\sqrt{13}+\sqrt5$.
At the positive minimizer both integrals in~\eqref{eq:integrals}
are strictly positive, and the jump in the diagonal derivative is $4/3$.
\end{proposition}

\begin{proof}
Here $a=\alpha$, $b=\beta$ and
\[
 P=x^6-7x^4+11x^2-1-2(\alpha+\beta)(x^3-3x)+4\alpha\beta.
\]
The kernel polynomials are
\[
 Z=1+7t^2+11t^4+t^6,\qquad
 X=Y=2t^3(1+3t^2),\qquad W=4t^6.
\]
Since $ZW-X^2=4t^8(1+t^2)^2>0$,
\[
 C-Z^2=X^2(\alpha^2+\beta^2)
          +2(X^2-ZW)\alpha\beta+W^2\alpha^2\beta^2
\]
is positive whenever $\alpha\beta\leq0$ and
$(\alpha,\beta)\ne(0,0)$. For positive product,
\[
 C(\alpha,\beta)-C(\sqrt{\alpha\beta},\sqrt{\alpha\beta})
      =X^2(\alpha-\beta)^2.
\]
Using simultaneous negation, every possible minimizer therefore
reduces to $(c,c)$, $0\leq c\leq1$.

On the diagonal,
\begin{equation}\label{eq:331-factors}
 P=(x^3-x^2-3x+1-2c)(x^3+x^2-3x-1-2c).
\end{equation}
Both cubics are real-rooted. For $0\leq c<1/2$, the first has one
negative root $-U$, and the second has one positive root $v$.
The products and traces of the roots give these sign counts.
Writing $h(z)=z^3+z^2-3z$, we have
\[
 U,v>1,\qquad h(U)=1-2c,\quad h(v)=1+2c,\qquad
 \E(c,c)=2+2(U+v).
\]
For $0<c<1/2$, $v>U$ and $h'$ is strictly increasing on $[1,\infty)$,
so
\[
 \frac{d}{dc}\E(c,c)
       =4\left(\frac1{h'(v)}-\frac1{h'(U)}\right)<0.
\]
For $1/2<c\leq1$, each cubic has one positive root, respectively
$w>2$ and $v>1$. Three positive roots are excluded by the negative
pairwise-root sum in the first cubic and the negative trace in
the second. Hence
\[
 \E(c,c)=2(w+v),\qquad
 \frac{d}{dc}\E(c,c)=\frac4{3w^2-2w-3}
                         +\frac4{3v^2+2v-3}>0.
\]
Continuity handles the endpoints and $c=1/2$. This proves the
asserted full-square minimizer set, including exclusion of the origin.
At $c=1/2$,
\[
 P=x(x+2)(x^2-x-3)(x^2-x-1),
\]
which gives the energy value already computed in
\cite[Proposition 4.9]{HuangPereyra2026}.
Each real-part pair has four cycle-product lifts, giving eight
switching classes.

The smooth part has diagonal derivative
\[
 G=-\frac25-\frac{2\sqrt{13}}{39}+\frac{2\sqrt5}{5},
       \qquad 0<G<\frac23 .
\]
The inequality follows, for example, from
$11/5<\sqrt5<9/4$ and $18/5<\sqrt{13}<37/10$.
Differentiating the extreme-root formulas above gives left and
right derivatives $G-2/3$ and $G+2/3$, respectively.
Their jump is $4/3$, as~\eqref{eq:zero-slope} predicts with $H=3$.
Symmetry gives $I_\alpha=I_\beta=\pi G/4$, so the strict
test~\eqref{eq:conditions} holds.
Directly, each stationarity numerator at this point is
\[
 -\tfrac12W(Z-W/4)+X^2
       =2t^6(1+5t^2+7t^4)>0.
\]
After division by $t^2C$ it is $O(t^4)$ at zero and $O(t^{-2})$
at infinity. Thus both integrals are positive at this interior
global minimizer, contradicting the unqualified zero equations
in~\cite[Theorem 4.7, equation (117)]{HuangPereyra2026}.
\end{proof}

\begin{proposition}\label{prop:335}
For $D_{3,3,5}$, at $(1/2,1/2)$ the directional energy derivative
along $(-1,-1)$ lies strictly between $-8/1000$ and $-7/1000$.
Consequently this singular point, and its simultaneous-negation
companion, are not local minima.
\end{proposition}

\begin{proof}
Let $P_c$ denote differentiation of $P(x;c,c)$ with respect to $c$,
evaluated at $c=1/2$, and set $Q(x)=P(x;1/2,1/2)$ for this proof.
Equation~\eqref{eq:matching} gives
\begin{align}
 Q(x)&=x(x+2)(x^4-x^3-5x^2+3x+5)
                    (x^4-x^3-3x^2+x+1),\label{eq:335-Q}\\
 P_c(x)&=-4(x^7-7x^5-x^4+13x^3+3x^2-5x-1).
                    \label{eq:335-Pc}
\end{align}
All ten roots of $Q$ are simple. Define the nonzero-root contribution
\[
 G_{335}=\sum_{\substack{Q(z)=0\\z\ne0}}
               -\sgn(z)\frac{P_c(z)}{Q'(z)}.
\]
The rational certificate in Appendix~\ref{app:certificate} proves
\[
 \frac{407}{1000}<G_{335}<\frac{408}{1000}.
\]
Implicit differentiation identifies this as the smooth diagonal
derivative. Here $H=5$ and the zero-root diagonal derivative is
$-2/5$. Therefore~\eqref{eq:expansion} gives the directional
energy derivative $2/5-G_{335}$, proving the bounds.
\end{proof}

\section{Scope and reproducibility}

Theorems~\ref{thm:singular} and~\ref{thm:product} give at most one
positive interior singular candidate for each length triple and
exclude all such candidates under~\eqref{eq:length-bound}.
They do not determine which candidates are global minima.
Equality in the normal bound, full-square sufficiency at singular
boundary endpoints, and the global minimum with all switching
classes for arbitrary odd $(r,s,\ell)$ remain unresolved here.
The $D_{3,3,5}$ certificate demonstrates why constrained minimality
and a zero-eigenvalue cusp alone do not suffice.

Symbolic checks used Python and SymPy 1.14.0~\cite{Meurer2017}.
They verified the coefficient identities on the 125 triples
$r,s\in\{3,5,7,9,11\}$, $\ell\in\{1,3,5,7,9\}$, along with
selected Hermitian characteristic polynomials and kernel identities.
These finite checks supplement the proofs for all lengths.
The $D_{3,3,5}$ certificate was separately recomputed by exact
rational Sturm sequences and midpoint Taylor enclosures; the
root-counting principle is the classical Sturm criterion
\cite[Section 1]{GonzalezEtAl1989}.
Floating-point quadrature, eigenvalue derivatives and energy
increments were diagnostic comparisons only.

\paragraph{Computational and AI assistance.}
OpenAI Codex assisted with proof development, symbolic and rational
verification programs, literature searches and manuscript drafting.
A separate model context reconstructed the mathematical arguments
and recomputed the decisive finite certificate. This constituted
an additional computational review pass, not external peer review.
The proofs and exact inequalities stated above are the basis of the
mathematical conclusions.

\appendix
\section{A rational certificate for the descent test}\label{app:certificate}

For the polynomials~\eqref{eq:335-Q}--\eqref{eq:335-Pc}, Table~\ref{tab:335}
lists isolating intervals for all nine nonzero roots. The remaining
root is exactly $0$. The last column encloses
$10^6[-\sgn(z)P_c(z)/Q'(z)]$; its endpoints are integers, so no
decimal approximation enters the certificate.

\begin{table}[htbp]
\centering
\small
\begin{tabular}{rrr}
\toprule
Lower endpoint & Upper endpoint & Scaled derivative enclosure\\
\midrule
$-2$ & $-2$ & $[-181819,-181818]$ \\
$-6062528/3328017$ & $-6253391/3432791$ & $[-351836,-351835]$ \\
$-1725791/1273013$ & $-10783276/7954179$ & $[-253297,-253296]$ \\
$-5627533/6650601$ & $-1669473/1972978$ & $[192844,192845]$ \\
$-6778623/14203208$ & $-977260/2047647$ & $[597764,597765]$ \\
$1273013/1725791$ & $7954179/10783276$ & $[-59616,-59615]$ \\
$2600452/1747381$ & $10436553/7012871$ & $[68791,68792]$ \\
$24441443/11664923$ & $2047647/977260$ & $[222264,222265]$ \\
$23045956/10573319$ & $22633443/10384061$ & $[172217,172218]$ \\
\midrule
\multicolumn{2}{r}{Sum of scaled derivative bounds} & $[407312,407321]$\\
\bottomrule
\end{tabular}
\caption{Exact nonzero-root intervals and derivative enclosures for $D_{3,3,5}$.}
\label{tab:335}
\end{table}

Here are explicit instructions to verify every entry. Form
$q_0=Q$, $q_1=Q'$ and
$q_{j+1}=-\operatorname{rem}(q_{j-1},q_j)$ over $\mathbb Q$ until
the remainder is zero. Positive rescaling of any $q_j$ is harmless.
The last nonzero polynomial is constant, so $\gcd(Q,Q')=1$.
Let $V(x)$ be the number of sign changes in $(q_j(x))$, with zero
entries removed. For each nonsingleton interval $[a,b]$ in the
table, $V(a)-V(b)=1$, and the endpoints are not roots.
The singleton $-2$ and the root $0$ are checked by direct
substitution, with $Q'$ nonzero. These ten disjoint root enclosures
exhaust the degree of $Q$.

To enclose a polynomial $R$ on $[a,b]$, put
$m=(a+b)/2$, $\rho=(b-a)/2$ and expand
$R(m+y)=\sum_{j=0}^d r_jy^j$ by rational arithmetic. Then
\[
 R([a,b])\subseteq
 \left[r_0-\sum_{j=1}^d|r_j|\rho^j,\
       r_0+\sum_{j=1}^d|r_j|\rho^j\right].
\]
Apply this to $P_c$ and $Q'$. Every denominator enclosure excludes
zero; divide the two intervals and multiply by $-\sgn(z)$.
Outward rounding after multiplication by $10^6$ gives the integer
endpoints in Table~\ref{tab:335}. Summing the lower and upper bounds
and dividing by $10^6$ proves the displayed bound for $G_{335}$.
The zero-root derivative is checked directly as
$-P_c(0)/Q'(0)=-2/5$.

\bibliographystyle{plain}
\bibliography{references}
\end{document}
